%% ============================================================
%% Lecture 7: Heteroskedasticity
%% Intermediate Econometrics — SUFE
%% Wooldridge, Introductory Econometrics, 8th edition, Chapter 8
%% ============================================================
\documentclass[aspectratio=169, 12pt]{beamer}
% ==============================================================================
% header.tex — LaTeX Preamble for Intermediate Econometrics
% Shandong University of Finance and Economics (SUFE)
% Textbook: Wooldridge, Introductory Econometrics, 8th edition
% ==============================================================================
%
% USAGE (from Slides/ directory):
%   % ==============================================================================
% header.tex — LaTeX Preamble for Intermediate Econometrics
% Shandong University of Finance and Economics (SUFE)
% Textbook: Wooldridge, Introductory Econometrics, 8th edition
% ==============================================================================
%
% USAGE (from Slides/ directory):
%   % ==============================================================================
% header.tex — LaTeX Preamble for Intermediate Econometrics
% Shandong University of Finance and Economics (SUFE)
% Textbook: Wooldridge, Introductory Econometrics, 8th edition
% ==============================================================================
%
% USAGE (from Slides/ directory):
%   \input{header}          % in Beamer .tex files (TEXINPUTS=../Preambles:)
%
% COMPILE (3-pass XeLaTeX):
%   cd Slides
%   TEXINPUTS=../Preambles:$TEXINPUTS xelatex -interaction=nonstopmode file.tex
%   BIBINPUTS=..:$BIBINPUTS bibtex file
%   TEXINPUTS=../Preambles:$TEXINPUTS xelatex -interaction=nonstopmode file.tex
%   TEXINPUTS=../Preambles:$TEXINPUTS xelatex -interaction=nonstopmode file.tex
% ==============================================================================


% --- Essential packages -------------------------------------------------------
\usepackage{fontspec}           % XeLaTeX font support (must load early)
\usepackage{amsmath}            % Math environments (align, equation, etc.)
\usepackage{amssymb}            % Math symbols (mathbb, etc.)
\usepackage{bm}                 % Bold math (\bm{})
\usepackage{mathtools}          % Extended math tools (coloneqq, etc.)


% --- Table packages -----------------------------------------------------------
\usepackage{booktabs}           % Professional tables (\toprule, \midrule, \bottomrule)
\usepackage{threeparttable}     % Table notes (tablenotes environment)
\usepackage{siunitx}            % Number alignment in tables (S column type)
\sisetup{
  table-format = 1.4,
  table-number-alignment = center,
  detect-weight = true,
  detect-inline-weight = math
}


% --- Figure and graphics packages --------------------------------------------
\usepackage{graphicx}           % Include figures (\includegraphics)
\usepackage{subcaption}         % Subfigures (subfigure environment)
\usepackage{tikz}               % TikZ diagrams
\usepackage{pgfplots}           % Data plots
\pgfplotsset{compat=1.18}
\usetikzlibrary{arrows.meta, positioning, shapes.geometric,
                decorations.pathreplacing, calc, patterns}


% --- Color and boxes ----------------------------------------------------------
\usepackage{xcolor}             % Extended color support
\usepackage{tcolorbox}          % Custom box environments
\tcbuselibrary{skins, breakable, theorems}


% --- Other utility packages ---------------------------------------------------
\usepackage{hyperref}           % Hyperlinks
\usepackage{natbib}             % Bibliography (\citet, \citep, \citealt)
\usepackage{caption}            % Figure/table captions
\usepackage{appendixnumberbeamer} % Appendix frame numbering


% ==============================================================================
% SUFE INSTITUTIONAL COLORS
% ==============================================================================
% Update hex values to match official SUFE brand colors if available.

\definecolor{SUFEBlue}{HTML}{012169}     % Primary: deep navy
\definecolor{SUFEGold}{HTML}{B9975B}     % Secondary: warm gold
\definecolor{SUFEYellow}{HTML}{F2A900}   % Highlight: bright gold/yellow
\definecolor{SUFELight}{HTML}{E8EDF5}    % Light background: pale blue-gray
\definecolor{SUFEDark}{HTML}{1a1a1a}     % Body text: near-black
\definecolor{positive}{HTML}{2E7D32}     % Positive/correct: dark green
\definecolor{negative}{HTML}{C62828}     % Negative/incorrect: dark red


% ==============================================================================
% BEAMER THEME & COLOR SETUP
% ==============================================================================

\usetheme{default}
\useinnertheme{default}
\useoutertheme{default}
\usecolortheme{default}

% Frame title
\setbeamercolor{frametitle}{fg=SUFEBlue, bg=white}
\setbeamerfont{frametitle}{size=\large, series=\bfseries}
\setbeamertemplate{frametitle}[default][left]
\addtobeamertemplate{frametitle}{\vspace{0.05em}}{%
  \vspace{0.05em}{\color{SUFEGold}\rule{\textwidth}{0.5pt}}\vspace{-0.1em}}

% Title page
\setbeamercolor{title}{fg=SUFEBlue}
\setbeamercolor{subtitle}{fg=SUFEGold}
\setbeamercolor{author}{fg=SUFEBlue}
\setbeamercolor{institute}{fg=SUFEDark}
\setbeamercolor{date}{fg=SUFEGold}

% Structure elements
\setbeamercolor{structure}{fg=SUFEBlue}
\setbeamercolor{item}{fg=SUFEBlue}
\setbeamercolor{subitem}{fg=SUFEGold}
\setbeamercolor{subsubitem}{fg=SUFEGold}
\setbeamerfont{itemize/enumerate body}{size=\normalsize}
\setbeamerfont{itemize/enumerate subbody}{size=\small}

% Block environments
\setbeamercolor{block title}{fg=white, bg=SUFEBlue}
\setbeamercolor{block body}{fg=SUFEDark, bg=SUFELight}
\setbeamercolor{block title alerted}{fg=white, bg=red!70!black}
\setbeamercolor{block body alerted}{fg=SUFEDark, bg=red!5}
\setbeamercolor{block title example}{fg=white, bg=SUFEGold!80!black}
\setbeamercolor{block body example}{fg=SUFEDark, bg=SUFEGold!8}

% Remove navigation symbols
\setbeamertemplate{navigation symbols}{}

% Footline: author | short title | slide N/N
\setbeamertemplate{footline}{%
  \leavevmode%
  \hbox{%
    \begin{beamercolorbox}[wd=.30\paperwidth,ht=2.25ex,dp=1ex,left,leftskip=2mm]{author in head/foot}%
      \usebeamerfont{author in head/foot}\insertshortauthor
    \end{beamercolorbox}%
    \begin{beamercolorbox}[wd=.40\paperwidth,ht=2.25ex,dp=1ex,center]{title in head/foot}%
      \usebeamerfont{title in head/foot}\insertshorttitle
    \end{beamercolorbox}%
    \begin{beamercolorbox}[wd=.30\paperwidth,ht=2.25ex,dp=1ex,right,rightskip=2mm]{date in head/foot}%
      \usebeamerfont{date in head/foot}%
      \insertframenumber{} / \inserttotalframenumber
    \end{beamercolorbox}}%
  \vskip0pt%
}
\setbeamercolor{author in head/foot}{fg=white, bg=SUFEBlue}
\setbeamercolor{title in head/foot}{fg=white, bg=SUFEBlue!80!black}
\setbeamercolor{date in head/foot}{fg=white, bg=SUFEBlue}

% Section separator slides (large centered section title)
\AtBeginSection[]{
  \begin{frame}
    \centering
    \vfill
    {\Large\bfseries\color{SUFEBlue}\insertsectionhead}
    \vfill
    {\color{SUFEGold}\rule{0.5\textwidth}{1pt}}
    \vfill
  \end{frame}
}


% ==============================================================================
% FONT SETUP (XeLaTeX)
% ==============================================================================
% Using Windows system fonts (Times New Roman / Arial / Consolas).
% On macOS, swap for: TeX Gyre Termes / TeX Gyre Heros / TeX Gyre Cursor
%   or: Times New Roman / Helvetica / Courier New

\setmainfont{Times New Roman}
\setsansfont{Arial}
\setmonofont{Consolas}


% ==============================================================================
% CUSTOM BOX ENVIRONMENTS (tcolorbox)
% ==============================================================================
% Quarto CSS equivalents: .keybox, .highlightbox, .methodbox,
%                         .assumptionbox, .resultbox

% keybox: key definitions, stated assumptions
\newtcolorbox{keybox}{
  enhanced,
  colback  = SUFEGold!10!white,
  colframe = SUFEGold,
  leftrule = 4pt, rightrule = 0pt, toprule = 0pt, bottomrule = 0pt,
  arc = 0pt, outer arc = 0pt,
  boxsep = 0pt, left = 8pt, right = 6pt, top = 4pt, bottom = 4pt,
}

% highlightbox: important theorems, main results
\newtcolorbox{highlightbox}{
  enhanced,
  colback  = SUFEYellow!8!white,
  colframe = SUFEYellow!80!SUFEGold,
  leftrule = 4pt, rightrule = 0pt, toprule = 0pt, bottomrule = 0pt,
  arc = 0pt, outer arc = 0pt,
  boxsep = 0pt, left = 8pt, right = 6pt, top = 4pt, bottom = 4pt,
}

% definitionbox{Title}: formal definitions — OLS, IV, GMM, etc.
\newtcolorbox{definitionbox}[1]{
  enhanced,
  title    = #1,
  colback  = SUFEBlue!4!white,
  colframe = SUFEBlue,
  colbacktitle = SUFEBlue,
  coltitle = white,
  fonttitle= \bfseries\small,
  arc = 4pt, outer arc = 4pt,
  boxsep = 0pt, left = 8pt, right = 6pt, top = 4pt, bottom = 4pt,
  toptitle = 3pt, bottomtitle = 3pt,
}

% examplebox{Title}: empirical examples, Wooldridge data applications
\newtcolorbox{examplebox}[1]{
  enhanced,
  title    = #1,
  colback  = SUFELight,
  colframe = SUFEGold!70!white,
  colbacktitle = SUFEGold!35!white,
  coltitle = SUFEBlue,
  fonttitle= \bfseries\small,
  arc = 4pt, outer arc = 4pt,
  boxsep = 0pt, left = 8pt, right = 6pt, top = 4pt, bottom = 4pt,
  toptitle = 3pt, bottomtitle = 3pt,
}

% assumptionbox: numbered OLS assumptions (MLR.1–MLR.6)
\newtcolorbox{assumptionbox}{
  enhanced,
  colback  = SUFEGold!8!white,
  colframe = SUFEGold,
  arc = 4pt, outer arc = 4pt,
  boxsep = 0pt, left = 8pt, right = 6pt, top = 4pt, bottom = 4pt,
}


% ==============================================================================
% STANDARD MATH MACROS
% ==============================================================================

% --- Operators ----------------------------------------------------------------
\DeclareMathOperator{\E}{\mathbb{E}}       % Expectation
\DeclareMathOperator{\Var}{Var}             % Variance
\DeclareMathOperator{\Cov}{Cov}             % Covariance
\DeclareMathOperator{\Corr}{Corr}           % Correlation
\DeclareMathOperator{\plim}{plim}           % Probability limit
\DeclareMathOperator{\rank}{rank}           % Matrix rank
\DeclareMathOperator{\tr}{tr}               % Matrix trace
\DeclareMathOperator{\se}{se}               % Standard error

% --- Common shortcuts ---------------------------------------------------------
\newcommand{\bhat}{\hat{\beta}}             % OLS point estimator
\newcommand{\bhatt}{\tilde{\beta}}          % Alternative estimator
\newcommand{\eps}{\varepsilon}              % Error term
\newcommand{\epsh}{\hat{\varepsilon}}       % OLS residual
\newcommand{\uh}{\hat{u}}                   % Residual (alternative)
\newcommand{\yh}{\hat{y}}                   % Fitted value

% Matrices and vectors
\newcommand{\Xb}{\mathbf{X}}               % Design matrix
\newcommand{\yb}{\mathbf{y}}               % Outcome vector
\newcommand{\betab}{\boldsymbol{\beta}}    % Parameter vector
\newcommand{\epsb}{\boldsymbol{\varepsilon}} % Error vector

% Goodness-of-fit
\newcommand{\SSR}{\text{SSR}}              % Sum of squared residuals
\newcommand{\SSE}{\text{SSE}}              % Explained sum of squares
\newcommand{\SST}{\text{SST}}              % Total sum of squares
\newcommand{\Rsq}{R^2}                     % R-squared
\newcommand{\aRsq}{\bar{R}^2}             % Adjusted R-squared

% Asymptotic notation
\newcommand{\avar}{\text{Avar}}            % Asymptotic variance
\newcommand{\pto}{\overset{p}{\to}}        % Convergence in probability
\newcommand{\dto}{\overset{d}{\to}}        % Convergence in distribution
\newcommand{\iid}{\overset{\text{iid}}{\sim}}  % i.i.d. distributed
\newcommand{\asim}{\overset{a}{\sim}}      % Approximately distributed

% Econometric estimators
\newcommand{\OLS}{\text{OLS}}
\newcommand{\IV}{\text{IV}}
\newcommand{\TSLS}{\text{2SLS}}
\newcommand{\GMM}{\text{GMM}}
\newcommand{\FE}{\text{FE}}
\newcommand{\RE}{\text{RE}}
\newcommand{\DID}{\text{DiD}}
\newcommand{\RD}{\text{RD}}

% Probability
\newcommand{\prob}{\mathbb{P}}             % Probability
\newcommand{\indep}{\perp\!\!\!\perp}      % Statistical independence

% Distributions
\newcommand{\Normal}{\mathcal{N}}          % Normal distribution
% Usage: X \sim \Normal(\mu, \sigma^2)

% Absolute value and norm
\newcommand{\abs}[1]{\left\lvert #1 \right\rvert}
\newcommand{\norm}[1]{\left\lVert #1 \right\rVert}


% ==============================================================================
% HYPERREF SETUP
% ==============================================================================

\hypersetup{
  colorlinks = true,
  linkcolor  = SUFEBlue,
  citecolor  = SUFEGold,
  urlcolor   = SUFEBlue
}


% ==============================================================================
% BIBLIOGRAPHY SETUP
% ==============================================================================

\bibliographystyle{aer}   % American Economic Review style
% Usage: \bibliography{../Bibliography_base}  (from Slides/ directory)
% In-text: \citet{Wooldridge2020}, \citep{White1980}
          % in Beamer .tex files (TEXINPUTS=../Preambles:)
%
% COMPILE (3-pass XeLaTeX):
%   cd Slides
%   TEXINPUTS=../Preambles:$TEXINPUTS xelatex -interaction=nonstopmode file.tex
%   BIBINPUTS=..:$BIBINPUTS bibtex file
%   TEXINPUTS=../Preambles:$TEXINPUTS xelatex -interaction=nonstopmode file.tex
%   TEXINPUTS=../Preambles:$TEXINPUTS xelatex -interaction=nonstopmode file.tex
% ==============================================================================


% --- Essential packages -------------------------------------------------------
\usepackage{fontspec}           % XeLaTeX font support (must load early)
\usepackage{amsmath}            % Math environments (align, equation, etc.)
\usepackage{amssymb}            % Math symbols (mathbb, etc.)
\usepackage{bm}                 % Bold math (\bm{})
\usepackage{mathtools}          % Extended math tools (coloneqq, etc.)


% --- Table packages -----------------------------------------------------------
\usepackage{booktabs}           % Professional tables (\toprule, \midrule, \bottomrule)
\usepackage{threeparttable}     % Table notes (tablenotes environment)
\usepackage{siunitx}            % Number alignment in tables (S column type)
\sisetup{
  table-format = 1.4,
  table-number-alignment = center,
  detect-weight = true,
  detect-inline-weight = math
}


% --- Figure and graphics packages --------------------------------------------
\usepackage{graphicx}           % Include figures (\includegraphics)
\usepackage{subcaption}         % Subfigures (subfigure environment)
\usepackage{tikz}               % TikZ diagrams
\usepackage{pgfplots}           % Data plots
\pgfplotsset{compat=1.18}
\usetikzlibrary{arrows.meta, positioning, shapes.geometric,
                decorations.pathreplacing, calc, patterns}


% --- Color and boxes ----------------------------------------------------------
\usepackage{xcolor}             % Extended color support
\usepackage{tcolorbox}          % Custom box environments
\tcbuselibrary{skins, breakable, theorems}


% --- Other utility packages ---------------------------------------------------
\usepackage{hyperref}           % Hyperlinks
\usepackage{natbib}             % Bibliography (\citet, \citep, \citealt)
\usepackage{caption}            % Figure/table captions
\usepackage{appendixnumberbeamer} % Appendix frame numbering


% ==============================================================================
% SUFE INSTITUTIONAL COLORS
% ==============================================================================
% Update hex values to match official SUFE brand colors if available.

\definecolor{SUFEBlue}{HTML}{012169}     % Primary: deep navy
\definecolor{SUFEGold}{HTML}{B9975B}     % Secondary: warm gold
\definecolor{SUFEYellow}{HTML}{F2A900}   % Highlight: bright gold/yellow
\definecolor{SUFELight}{HTML}{E8EDF5}    % Light background: pale blue-gray
\definecolor{SUFEDark}{HTML}{1a1a1a}     % Body text: near-black
\definecolor{positive}{HTML}{2E7D32}     % Positive/correct: dark green
\definecolor{negative}{HTML}{C62828}     % Negative/incorrect: dark red


% ==============================================================================
% BEAMER THEME & COLOR SETUP
% ==============================================================================

\usetheme{default}
\useinnertheme{default}
\useoutertheme{default}
\usecolortheme{default}

% Frame title
\setbeamercolor{frametitle}{fg=SUFEBlue, bg=white}
\setbeamerfont{frametitle}{size=\large, series=\bfseries}
\setbeamertemplate{frametitle}[default][left]
\addtobeamertemplate{frametitle}{\vspace{0.05em}}{%
  \vspace{0.05em}{\color{SUFEGold}\rule{\textwidth}{0.5pt}}\vspace{-0.1em}}

% Title page
\setbeamercolor{title}{fg=SUFEBlue}
\setbeamercolor{subtitle}{fg=SUFEGold}
\setbeamercolor{author}{fg=SUFEBlue}
\setbeamercolor{institute}{fg=SUFEDark}
\setbeamercolor{date}{fg=SUFEGold}

% Structure elements
\setbeamercolor{structure}{fg=SUFEBlue}
\setbeamercolor{item}{fg=SUFEBlue}
\setbeamercolor{subitem}{fg=SUFEGold}
\setbeamercolor{subsubitem}{fg=SUFEGold}
\setbeamerfont{itemize/enumerate body}{size=\normalsize}
\setbeamerfont{itemize/enumerate subbody}{size=\small}

% Block environments
\setbeamercolor{block title}{fg=white, bg=SUFEBlue}
\setbeamercolor{block body}{fg=SUFEDark, bg=SUFELight}
\setbeamercolor{block title alerted}{fg=white, bg=red!70!black}
\setbeamercolor{block body alerted}{fg=SUFEDark, bg=red!5}
\setbeamercolor{block title example}{fg=white, bg=SUFEGold!80!black}
\setbeamercolor{block body example}{fg=SUFEDark, bg=SUFEGold!8}

% Remove navigation symbols
\setbeamertemplate{navigation symbols}{}

% Footline: author | short title | slide N/N
\setbeamertemplate{footline}{%
  \leavevmode%
  \hbox{%
    \begin{beamercolorbox}[wd=.30\paperwidth,ht=2.25ex,dp=1ex,left,leftskip=2mm]{author in head/foot}%
      \usebeamerfont{author in head/foot}\insertshortauthor
    \end{beamercolorbox}%
    \begin{beamercolorbox}[wd=.40\paperwidth,ht=2.25ex,dp=1ex,center]{title in head/foot}%
      \usebeamerfont{title in head/foot}\insertshorttitle
    \end{beamercolorbox}%
    \begin{beamercolorbox}[wd=.30\paperwidth,ht=2.25ex,dp=1ex,right,rightskip=2mm]{date in head/foot}%
      \usebeamerfont{date in head/foot}%
      \insertframenumber{} / \inserttotalframenumber
    \end{beamercolorbox}}%
  \vskip0pt%
}
\setbeamercolor{author in head/foot}{fg=white, bg=SUFEBlue}
\setbeamercolor{title in head/foot}{fg=white, bg=SUFEBlue!80!black}
\setbeamercolor{date in head/foot}{fg=white, bg=SUFEBlue}

% Section separator slides (large centered section title)
\AtBeginSection[]{
  \begin{frame}
    \centering
    \vfill
    {\Large\bfseries\color{SUFEBlue}\insertsectionhead}
    \vfill
    {\color{SUFEGold}\rule{0.5\textwidth}{1pt}}
    \vfill
  \end{frame}
}


% ==============================================================================
% FONT SETUP (XeLaTeX)
% ==============================================================================
% Using Windows system fonts (Times New Roman / Arial / Consolas).
% On macOS, swap for: TeX Gyre Termes / TeX Gyre Heros / TeX Gyre Cursor
%   or: Times New Roman / Helvetica / Courier New

\setmainfont{Times New Roman}
\setsansfont{Arial}
\setmonofont{Consolas}


% ==============================================================================
% CUSTOM BOX ENVIRONMENTS (tcolorbox)
% ==============================================================================
% Quarto CSS equivalents: .keybox, .highlightbox, .methodbox,
%                         .assumptionbox, .resultbox

% keybox: key definitions, stated assumptions
\newtcolorbox{keybox}{
  enhanced,
  colback  = SUFEGold!10!white,
  colframe = SUFEGold,
  leftrule = 4pt, rightrule = 0pt, toprule = 0pt, bottomrule = 0pt,
  arc = 0pt, outer arc = 0pt,
  boxsep = 0pt, left = 8pt, right = 6pt, top = 4pt, bottom = 4pt,
}

% highlightbox: important theorems, main results
\newtcolorbox{highlightbox}{
  enhanced,
  colback  = SUFEYellow!8!white,
  colframe = SUFEYellow!80!SUFEGold,
  leftrule = 4pt, rightrule = 0pt, toprule = 0pt, bottomrule = 0pt,
  arc = 0pt, outer arc = 0pt,
  boxsep = 0pt, left = 8pt, right = 6pt, top = 4pt, bottom = 4pt,
}

% definitionbox{Title}: formal definitions — OLS, IV, GMM, etc.
\newtcolorbox{definitionbox}[1]{
  enhanced,
  title    = #1,
  colback  = SUFEBlue!4!white,
  colframe = SUFEBlue,
  colbacktitle = SUFEBlue,
  coltitle = white,
  fonttitle= \bfseries\small,
  arc = 4pt, outer arc = 4pt,
  boxsep = 0pt, left = 8pt, right = 6pt, top = 4pt, bottom = 4pt,
  toptitle = 3pt, bottomtitle = 3pt,
}

% examplebox{Title}: empirical examples, Wooldridge data applications
\newtcolorbox{examplebox}[1]{
  enhanced,
  title    = #1,
  colback  = SUFELight,
  colframe = SUFEGold!70!white,
  colbacktitle = SUFEGold!35!white,
  coltitle = SUFEBlue,
  fonttitle= \bfseries\small,
  arc = 4pt, outer arc = 4pt,
  boxsep = 0pt, left = 8pt, right = 6pt, top = 4pt, bottom = 4pt,
  toptitle = 3pt, bottomtitle = 3pt,
}

% assumptionbox: numbered OLS assumptions (MLR.1–MLR.6)
\newtcolorbox{assumptionbox}{
  enhanced,
  colback  = SUFEGold!8!white,
  colframe = SUFEGold,
  arc = 4pt, outer arc = 4pt,
  boxsep = 0pt, left = 8pt, right = 6pt, top = 4pt, bottom = 4pt,
}


% ==============================================================================
% STANDARD MATH MACROS
% ==============================================================================

% --- Operators ----------------------------------------------------------------
\DeclareMathOperator{\E}{\mathbb{E}}       % Expectation
\DeclareMathOperator{\Var}{Var}             % Variance
\DeclareMathOperator{\Cov}{Cov}             % Covariance
\DeclareMathOperator{\Corr}{Corr}           % Correlation
\DeclareMathOperator{\plim}{plim}           % Probability limit
\DeclareMathOperator{\rank}{rank}           % Matrix rank
\DeclareMathOperator{\tr}{tr}               % Matrix trace
\DeclareMathOperator{\se}{se}               % Standard error

% --- Common shortcuts ---------------------------------------------------------
\newcommand{\bhat}{\hat{\beta}}             % OLS point estimator
\newcommand{\bhatt}{\tilde{\beta}}          % Alternative estimator
\newcommand{\eps}{\varepsilon}              % Error term
\newcommand{\epsh}{\hat{\varepsilon}}       % OLS residual
\newcommand{\uh}{\hat{u}}                   % Residual (alternative)
\newcommand{\yh}{\hat{y}}                   % Fitted value

% Matrices and vectors
\newcommand{\Xb}{\mathbf{X}}               % Design matrix
\newcommand{\yb}{\mathbf{y}}               % Outcome vector
\newcommand{\betab}{\boldsymbol{\beta}}    % Parameter vector
\newcommand{\epsb}{\boldsymbol{\varepsilon}} % Error vector

% Goodness-of-fit
\newcommand{\SSR}{\text{SSR}}              % Sum of squared residuals
\newcommand{\SSE}{\text{SSE}}              % Explained sum of squares
\newcommand{\SST}{\text{SST}}              % Total sum of squares
\newcommand{\Rsq}{R^2}                     % R-squared
\newcommand{\aRsq}{\bar{R}^2}             % Adjusted R-squared

% Asymptotic notation
\newcommand{\avar}{\text{Avar}}            % Asymptotic variance
\newcommand{\pto}{\overset{p}{\to}}        % Convergence in probability
\newcommand{\dto}{\overset{d}{\to}}        % Convergence in distribution
\newcommand{\iid}{\overset{\text{iid}}{\sim}}  % i.i.d. distributed
\newcommand{\asim}{\overset{a}{\sim}}      % Approximately distributed

% Econometric estimators
\newcommand{\OLS}{\text{OLS}}
\newcommand{\IV}{\text{IV}}
\newcommand{\TSLS}{\text{2SLS}}
\newcommand{\GMM}{\text{GMM}}
\newcommand{\FE}{\text{FE}}
\newcommand{\RE}{\text{RE}}
\newcommand{\DID}{\text{DiD}}
\newcommand{\RD}{\text{RD}}

% Probability
\newcommand{\prob}{\mathbb{P}}             % Probability
\newcommand{\indep}{\perp\!\!\!\perp}      % Statistical independence

% Distributions
\newcommand{\Normal}{\mathcal{N}}          % Normal distribution
% Usage: X \sim \Normal(\mu, \sigma^2)

% Absolute value and norm
\newcommand{\abs}[1]{\left\lvert #1 \right\rvert}
\newcommand{\norm}[1]{\left\lVert #1 \right\rVert}


% ==============================================================================
% HYPERREF SETUP
% ==============================================================================

\hypersetup{
  colorlinks = true,
  linkcolor  = SUFEBlue,
  citecolor  = SUFEGold,
  urlcolor   = SUFEBlue
}


% ==============================================================================
% BIBLIOGRAPHY SETUP
% ==============================================================================

\bibliographystyle{aer}   % American Economic Review style
% Usage: \bibliography{../Bibliography_base}  (from Slides/ directory)
% In-text: \citet{Wooldridge2020}, \citep{White1980}
          % in Beamer .tex files (TEXINPUTS=../Preambles:)
%
% COMPILE (3-pass XeLaTeX):
%   cd Slides
%   TEXINPUTS=../Preambles:$TEXINPUTS xelatex -interaction=nonstopmode file.tex
%   BIBINPUTS=..:$BIBINPUTS bibtex file
%   TEXINPUTS=../Preambles:$TEXINPUTS xelatex -interaction=nonstopmode file.tex
%   TEXINPUTS=../Preambles:$TEXINPUTS xelatex -interaction=nonstopmode file.tex
% ==============================================================================


% --- Essential packages -------------------------------------------------------
\usepackage{fontspec}           % XeLaTeX font support (must load early)
\usepackage{amsmath}            % Math environments (align, equation, etc.)
\usepackage{amssymb}            % Math symbols (mathbb, etc.)
\usepackage{bm}                 % Bold math (\bm{})
\usepackage{mathtools}          % Extended math tools (coloneqq, etc.)


% --- Table packages -----------------------------------------------------------
\usepackage{booktabs}           % Professional tables (\toprule, \midrule, \bottomrule)
\usepackage{threeparttable}     % Table notes (tablenotes environment)
\usepackage{siunitx}            % Number alignment in tables (S column type)
\sisetup{
  table-format = 1.4,
  table-number-alignment = center,
  detect-weight = true,
  detect-inline-weight = math
}


% --- Figure and graphics packages --------------------------------------------
\usepackage{graphicx}           % Include figures (\includegraphics)
\usepackage{subcaption}         % Subfigures (subfigure environment)
\usepackage{tikz}               % TikZ diagrams
\usepackage{pgfplots}           % Data plots
\pgfplotsset{compat=1.18}
\usetikzlibrary{arrows.meta, positioning, shapes.geometric,
                decorations.pathreplacing, calc, patterns}


% --- Color and boxes ----------------------------------------------------------
\usepackage{xcolor}             % Extended color support
\usepackage{tcolorbox}          % Custom box environments
\tcbuselibrary{skins, breakable, theorems}


% --- Other utility packages ---------------------------------------------------
\usepackage{hyperref}           % Hyperlinks
\usepackage{natbib}             % Bibliography (\citet, \citep, \citealt)
\usepackage{caption}            % Figure/table captions
\usepackage{appendixnumberbeamer} % Appendix frame numbering


% ==============================================================================
% SUFE INSTITUTIONAL COLORS
% ==============================================================================
% Update hex values to match official SUFE brand colors if available.

\definecolor{SUFEBlue}{HTML}{012169}     % Primary: deep navy
\definecolor{SUFEGold}{HTML}{B9975B}     % Secondary: warm gold
\definecolor{SUFEYellow}{HTML}{F2A900}   % Highlight: bright gold/yellow
\definecolor{SUFELight}{HTML}{E8EDF5}    % Light background: pale blue-gray
\definecolor{SUFEDark}{HTML}{1a1a1a}     % Body text: near-black
\definecolor{positive}{HTML}{2E7D32}     % Positive/correct: dark green
\definecolor{negative}{HTML}{C62828}     % Negative/incorrect: dark red


% ==============================================================================
% BEAMER THEME & COLOR SETUP
% ==============================================================================

\usetheme{default}
\useinnertheme{default}
\useoutertheme{default}
\usecolortheme{default}

% Frame title
\setbeamercolor{frametitle}{fg=SUFEBlue, bg=white}
\setbeamerfont{frametitle}{size=\large, series=\bfseries}
\setbeamertemplate{frametitle}[default][left]
\addtobeamertemplate{frametitle}{\vspace{0.05em}}{%
  \vspace{0.05em}{\color{SUFEGold}\rule{\textwidth}{0.5pt}}\vspace{-0.1em}}

% Title page
\setbeamercolor{title}{fg=SUFEBlue}
\setbeamercolor{subtitle}{fg=SUFEGold}
\setbeamercolor{author}{fg=SUFEBlue}
\setbeamercolor{institute}{fg=SUFEDark}
\setbeamercolor{date}{fg=SUFEGold}

% Structure elements
\setbeamercolor{structure}{fg=SUFEBlue}
\setbeamercolor{item}{fg=SUFEBlue}
\setbeamercolor{subitem}{fg=SUFEGold}
\setbeamercolor{subsubitem}{fg=SUFEGold}
\setbeamerfont{itemize/enumerate body}{size=\normalsize}
\setbeamerfont{itemize/enumerate subbody}{size=\small}

% Block environments
\setbeamercolor{block title}{fg=white, bg=SUFEBlue}
\setbeamercolor{block body}{fg=SUFEDark, bg=SUFELight}
\setbeamercolor{block title alerted}{fg=white, bg=red!70!black}
\setbeamercolor{block body alerted}{fg=SUFEDark, bg=red!5}
\setbeamercolor{block title example}{fg=white, bg=SUFEGold!80!black}
\setbeamercolor{block body example}{fg=SUFEDark, bg=SUFEGold!8}

% Remove navigation symbols
\setbeamertemplate{navigation symbols}{}

% Footline: author | short title | slide N/N
\setbeamertemplate{footline}{%
  \leavevmode%
  \hbox{%
    \begin{beamercolorbox}[wd=.30\paperwidth,ht=2.25ex,dp=1ex,left,leftskip=2mm]{author in head/foot}%
      \usebeamerfont{author in head/foot}\insertshortauthor
    \end{beamercolorbox}%
    \begin{beamercolorbox}[wd=.40\paperwidth,ht=2.25ex,dp=1ex,center]{title in head/foot}%
      \usebeamerfont{title in head/foot}\insertshorttitle
    \end{beamercolorbox}%
    \begin{beamercolorbox}[wd=.30\paperwidth,ht=2.25ex,dp=1ex,right,rightskip=2mm]{date in head/foot}%
      \usebeamerfont{date in head/foot}%
      \insertframenumber{} / \inserttotalframenumber
    \end{beamercolorbox}}%
  \vskip0pt%
}
\setbeamercolor{author in head/foot}{fg=white, bg=SUFEBlue}
\setbeamercolor{title in head/foot}{fg=white, bg=SUFEBlue!80!black}
\setbeamercolor{date in head/foot}{fg=white, bg=SUFEBlue}

% Section separator slides (large centered section title)
\AtBeginSection[]{
  \begin{frame}
    \centering
    \vfill
    {\Large\bfseries\color{SUFEBlue}\insertsectionhead}
    \vfill
    {\color{SUFEGold}\rule{0.5\textwidth}{1pt}}
    \vfill
  \end{frame}
}


% ==============================================================================
% FONT SETUP (XeLaTeX)
% ==============================================================================
% Using Windows system fonts (Times New Roman / Arial / Consolas).
% On macOS, swap for: TeX Gyre Termes / TeX Gyre Heros / TeX Gyre Cursor
%   or: Times New Roman / Helvetica / Courier New

\setmainfont{Times New Roman}
\setsansfont{Arial}
\setmonofont{Consolas}


% ==============================================================================
% CUSTOM BOX ENVIRONMENTS (tcolorbox)
% ==============================================================================
% Quarto CSS equivalents: .keybox, .highlightbox, .methodbox,
%                         .assumptionbox, .resultbox

% keybox: key definitions, stated assumptions
\newtcolorbox{keybox}{
  enhanced,
  colback  = SUFEGold!10!white,
  colframe = SUFEGold,
  leftrule = 4pt, rightrule = 0pt, toprule = 0pt, bottomrule = 0pt,
  arc = 0pt, outer arc = 0pt,
  boxsep = 0pt, left = 8pt, right = 6pt, top = 4pt, bottom = 4pt,
}

% highlightbox: important theorems, main results
\newtcolorbox{highlightbox}{
  enhanced,
  colback  = SUFEYellow!8!white,
  colframe = SUFEYellow!80!SUFEGold,
  leftrule = 4pt, rightrule = 0pt, toprule = 0pt, bottomrule = 0pt,
  arc = 0pt, outer arc = 0pt,
  boxsep = 0pt, left = 8pt, right = 6pt, top = 4pt, bottom = 4pt,
}

% definitionbox{Title}: formal definitions — OLS, IV, GMM, etc.
\newtcolorbox{definitionbox}[1]{
  enhanced,
  title    = #1,
  colback  = SUFEBlue!4!white,
  colframe = SUFEBlue,
  colbacktitle = SUFEBlue,
  coltitle = white,
  fonttitle= \bfseries\small,
  arc = 4pt, outer arc = 4pt,
  boxsep = 0pt, left = 8pt, right = 6pt, top = 4pt, bottom = 4pt,
  toptitle = 3pt, bottomtitle = 3pt,
}

% examplebox{Title}: empirical examples, Wooldridge data applications
\newtcolorbox{examplebox}[1]{
  enhanced,
  title    = #1,
  colback  = SUFELight,
  colframe = SUFEGold!70!white,
  colbacktitle = SUFEGold!35!white,
  coltitle = SUFEBlue,
  fonttitle= \bfseries\small,
  arc = 4pt, outer arc = 4pt,
  boxsep = 0pt, left = 8pt, right = 6pt, top = 4pt, bottom = 4pt,
  toptitle = 3pt, bottomtitle = 3pt,
}

% assumptionbox: numbered OLS assumptions (MLR.1–MLR.6)
\newtcolorbox{assumptionbox}{
  enhanced,
  colback  = SUFEGold!8!white,
  colframe = SUFEGold,
  arc = 4pt, outer arc = 4pt,
  boxsep = 0pt, left = 8pt, right = 6pt, top = 4pt, bottom = 4pt,
}


% ==============================================================================
% STANDARD MATH MACROS
% ==============================================================================

% --- Operators ----------------------------------------------------------------
\DeclareMathOperator{\E}{\mathbb{E}}       % Expectation
\DeclareMathOperator{\Var}{Var}             % Variance
\DeclareMathOperator{\Cov}{Cov}             % Covariance
\DeclareMathOperator{\Corr}{Corr}           % Correlation
\DeclareMathOperator{\plim}{plim}           % Probability limit
\DeclareMathOperator{\rank}{rank}           % Matrix rank
\DeclareMathOperator{\tr}{tr}               % Matrix trace
\DeclareMathOperator{\se}{se}               % Standard error

% --- Common shortcuts ---------------------------------------------------------
\newcommand{\bhat}{\hat{\beta}}             % OLS point estimator
\newcommand{\bhatt}{\tilde{\beta}}          % Alternative estimator
\newcommand{\eps}{\varepsilon}              % Error term
\newcommand{\epsh}{\hat{\varepsilon}}       % OLS residual
\newcommand{\uh}{\hat{u}}                   % Residual (alternative)
\newcommand{\yh}{\hat{y}}                   % Fitted value

% Matrices and vectors
\newcommand{\Xb}{\mathbf{X}}               % Design matrix
\newcommand{\yb}{\mathbf{y}}               % Outcome vector
\newcommand{\betab}{\boldsymbol{\beta}}    % Parameter vector
\newcommand{\epsb}{\boldsymbol{\varepsilon}} % Error vector

% Goodness-of-fit
\newcommand{\SSR}{\text{SSR}}              % Sum of squared residuals
\newcommand{\SSE}{\text{SSE}}              % Explained sum of squares
\newcommand{\SST}{\text{SST}}              % Total sum of squares
\newcommand{\Rsq}{R^2}                     % R-squared
\newcommand{\aRsq}{\bar{R}^2}             % Adjusted R-squared

% Asymptotic notation
\newcommand{\avar}{\text{Avar}}            % Asymptotic variance
\newcommand{\pto}{\overset{p}{\to}}        % Convergence in probability
\newcommand{\dto}{\overset{d}{\to}}        % Convergence in distribution
\newcommand{\iid}{\overset{\text{iid}}{\sim}}  % i.i.d. distributed
\newcommand{\asim}{\overset{a}{\sim}}      % Approximately distributed

% Econometric estimators
\newcommand{\OLS}{\text{OLS}}
\newcommand{\IV}{\text{IV}}
\newcommand{\TSLS}{\text{2SLS}}
\newcommand{\GMM}{\text{GMM}}
\newcommand{\FE}{\text{FE}}
\newcommand{\RE}{\text{RE}}
\newcommand{\DID}{\text{DiD}}
\newcommand{\RD}{\text{RD}}

% Probability
\newcommand{\prob}{\mathbb{P}}             % Probability
\newcommand{\indep}{\perp\!\!\!\perp}      % Statistical independence

% Distributions
\newcommand{\Normal}{\mathcal{N}}          % Normal distribution
% Usage: X \sim \Normal(\mu, \sigma^2)

% Absolute value and norm
\newcommand{\abs}[1]{\left\lvert #1 \right\rvert}
\newcommand{\norm}[1]{\left\lVert #1 \right\rVert}


% ==============================================================================
% HYPERREF SETUP
% ==============================================================================

\hypersetup{
  colorlinks = true,
  linkcolor  = SUFEBlue,
  citecolor  = SUFEGold,
  urlcolor   = SUFEBlue
}


% ==============================================================================
% BIBLIOGRAPHY SETUP
% ==============================================================================

\bibliographystyle{aer}   % American Economic Review style
% Usage: \bibliography{../Bibliography_base}  (from Slides/ directory)
% In-text: \citet{Wooldridge2020}, \citep{White1980}


\title[L7: Heteroskedasticity]{Heteroskedasticity}
\subtitle{Intermediate Econometrics}
\author[SUFE Econometrics]{Chengye Jia}
\institute{%
  Department of Economics \\
  Shandong University of Finance and Economics}
\date{Spring 2026 \\ \smallskip \scriptsize Wooldridge, Chapter 8}

%% ============================================================
\begin{document}
%% ============================================================

\begin{frame}[plain]
  \titlepage
\end{frame}

\begin{frame}[shrink=10]{Outline}
  \tableofcontents
\end{frame}

%% ============================================================
\section{Motivation and Definition (\S8-0)}
%% ============================================================

\begin{frame}{What Is Heteroskedasticity?}

  Recall \textbf{MLR.5 (Homoskedasticity)}:
  \[
    \Var(u \mid x_1, \ldots, x_k) = \sigma^2 \quad \text{(constant for all observations)}
  \]

  \smallskip

  \begin{definitionbox}{Heteroskedasticity}
    MLR.5 is \textbf{violated} when the error variance depends on the regressors:
    \[
      \Var(u_i \mid x_{i1}, \ldots, x_{ik}) = \sigma_i^2
    \]
    where $\sigma_i^2$ \textbf{varies} across observations $i = 1, \ldots, n$.
  \end{definitionbox}

  \smallskip

  \begin{keybox}
    \textbf{Why does this happen?}
    \begin{itemize}
      \item Richer households have \textit{more variable} consumption (income effect on spread)
      \item Larger firms have \textit{more variable} profits than small firms
      \item Binary outcomes: $\Var(y \mid x) = p(x)[1-p(x)]$ depends on $x$ by construction
    \end{itemize}
  \end{keybox}

\end{frame}

\begin{frame}{Visualising Heteroskedasticity}

  \begin{columns}[t]
    \column{0.50\textwidth}
    \textbf{Homoskedastic errors}

    \vspace{0.3cm}

    \begin{tikzpicture}[scale=0.88]
      \draw[->,line width=1pt] (0,0) -- (5.2,0) node[below] {$x$};
      \draw[->,line width=1pt] (0,-2) -- (0,2) node[left] {$\hat{u}$};
      \draw[dashed,gray] (0,0) -- (5,0);
      \foreach \x in {0.5,1.0,1.5,2.0,2.5,3.0,3.5,4.0,4.5,5.0}{
        \pgfmathsetmacro{\ya}{0.9*rand}
        \fill[SUFEBlue] (\x,\ya) circle (1.5pt);
      }
      \draw[dashed,SUFEGold] (0,-1.1) -- (5,-1.1);
      \draw[dashed,SUFEGold] (0, 1.1) -- (5, 1.1);
      \node[below,font=\small] at (2.5,-1.9) {Constant band};
    \end{tikzpicture}

    \column{0.47\textwidth}
    \textbf{Heteroskedastic errors}

    \vspace{0.3cm}

    \begin{tikzpicture}[scale=0.88]
      \draw[->,line width=1pt] (0,0) -- (5.2,0) node[below] {$x$};
      \draw[->,line width=1pt] (0,-2) -- (0,2) node[left] {$\hat{u}$};
      \draw[dashed,gray] (0,0) -- (5,0);
      \foreach \x/\s in {0.5/0.12,1.0/0.25,1.5/0.40,2.0/0.60,2.5/0.78,3.0/1.00,3.5/1.18,4.0/1.45,4.5/1.60,5.0/1.80}{
        \pgfmathsetmacro{\ya}{\s*rand}
        \fill[red!70] (\x,\ya) circle (1.5pt);
      }
      \draw[dashed,red!60] (0,-0.2) -- (5,-1.9);
      \draw[dashed,red!60] (0, 0.2) -- (5, 1.9);
      \node[below,font=\small] at (2.5,-1.9) {Fan-shaped spread};
    \end{tikzpicture}
  \end{columns}

  \smallskip

  \begin{keybox}
    Residual plots ($\hat{u}_i$ vs.\ $\hat{y}_i$ or $\hat{u}_i^2$ vs.\ $x_j$) are the first diagnostic. A funnel or fan shape signals heteroskedasticity.
  \end{keybox}

\end{frame}

%% ============================================================
\section{Consequences for OLS (\S8-1)}
%% ============================================================

\begin{frame}{Consequences: What Changes, What Stays the Same}

  \begin{columns}[t]
    \column{0.48\textwidth}
    \begin{highlightbox}
      \textbf{\textcolor{positive}{NOT affected:}}
      \begin{itemize}
        \item \textbf{Unbiasedness} of $\hat{\beta}_j$ (only MLR.1--4 needed)
        \item \textbf{Consistency} of $\hat{\beta}_j$
        \item Interpretation of $R^2$ and $\bar{R}^2$
      \end{itemize}
      \smallskip
      $\Rightarrow$ OLS point estimates are \textit{still useful}.
    \end{highlightbox}

    \column{0.48\textwidth}
    \begin{highlightbox}
      \textbf{\textcolor{negative}{IS affected:}}
      \begin{itemize}
        \item OLS \textbf{standard errors} are wrong
        \item $t$\textbf{-tests} and $F$\textbf{-tests} are invalid
        \item OLS is \textbf{no longer BLUE} (Gauss-Markov fails)
        \item \textbf{Efficiency} is lost
      \end{itemize}
      \smallskip
      $\Rightarrow$ Inference based on usual SEs is \textit{unreliable}.
    \end{highlightbox}
  \end{columns}

\end{frame}

\begin{frame}{Why OLS Standard Errors Are Wrong: The Proof}

  \textbf{True variance of $\hat{\beta}_1$ in simple regression} (derived from $\hat{\beta}_1 = \beta_1 + \sum_i w_i u_i$):
  \[
    \Var(\hat{\beta}_1) = \frac{\sum_{i=1}^n (x_i - \bar{x})^2 \sigma_i^2}{\SST_x^2} \quad \text{(eq.\ 8.2)}
  \]
  where $\SST_x = \sum_i(x_i-\bar{x})^2$.

  \smallskip

  \textbf{Under homoskedasticity} ($\sigma_i^2 = \sigma^2$ for all $i$), this simplifies to:
  \[
    \Var(\hat{\beta}_1) = \frac{\sigma^2 \sum_i(x_i-\bar{x})^2}{\SST_x^2} = \frac{\sigma^2}{\SST_x}
  \]

  \begin{keybox}
    OLS uses $\hat{\sigma}^2 = \SSR/(n-k-1)$ to estimate $\sigma^2$. If $\sigma_i^2 \neq \sigma^2$, then $\hat{\sigma}^2$ estimates a \textbf{weighted average} of $\sigma_i^2$, not the right quantity for eq.\ 8.2.

    \smallskip
    Direction of bias: \textit{unpredictable} without knowing how $\sigma_i^2$ relates to leverage $(x_i-\bar{x})^2$.
  \end{keybox}

\end{frame}

\begin{frame}{Gauss-Markov Fails: OLS Is No Longer BLUE}

  Recall the Gauss-Markov Theorem requires MLR.1--\textbf{5}. Under heteroskedasticity:

  \begin{highlightbox}
    There exist \textbf{other linear unbiased estimators} with \textit{smaller variance} than OLS.

    Specifically, \textbf{Weighted Least Squares} (WLS) with the correct weights achieves lower variance than OLS.
  \end{highlightbox}

  \smallskip

  \textbf{Intuition:}
  \begin{itemize}
    \item Observations with \textit{small} $\sigma_i^2$ are very informative — OLS under-uses them.
    \item Observations with \textit{large} $\sigma_i^2$ are noisy — OLS over-uses them.
    \item WLS corrects this by weighting inversely proportional to variance.
  \end{itemize}

  \smallskip

  \begin{keybox}
    Even though OLS estimates remain unbiased, an analyst using OLS under heteroskedasticity is \textbf{leaving efficiency on the table}. If we know the form of heteroskedasticity, WLS (§8-4) does better.
  \end{keybox}

\end{frame}

%% ============================================================
\section{Heteroskedasticity-Robust Inference (\S8-2)}
%% ============================================================

\begin{frame}{The White Robust Variance Estimator}

  \textbf{Key insight (White, 1980):} Replace the unknown $\sigma_i^2$ in eq.\ 8.2 with $\hat{u}_i^2$ (the squared OLS residual for observation $i$).

  \smallskip

  \begin{definitionbox}{Heteroskedasticity-Robust Variance Estimator (eq.\ 8.4)}
    \[
      \widehat{\Var}_{\text{rob}}(\hat{\beta}_j) = \frac{\displaystyle\sum_{i=1}^n \hat{r}_{ij}^2 \, \hat{u}_i^2}{\SSR_j^2}
    \]
    \begin{itemize}
      \item $\hat{r}_{ij}$ = residuals from regressing $x_j$ on \textit{all other regressors}
      \item $\SSR_j = \sum_i \hat{r}_{ij}^2$
      \item $\hat{u}_i$ = OLS residuals from the main regression
    \end{itemize}
  \end{definitionbox}

  \smallskip

  \begin{itemize}
    \item Also called: \textbf{White SE}, \textbf{Huber-White SE}, \textbf{HC standard error}.
    \item Under MLR.1--4, this estimator is \textbf{consistent} for any form of heteroskedasticity.
  \end{itemize}

\end{frame}

\begin{frame}{Why the Robust Estimator Is Consistent: Sketch}

  \textbf{We need:} A consistent estimator of $\sum_i \hat{r}_{ij}^2 \sigma_i^2 / \SSR_j^2$.

  \smallskip

  \textbf{Step 1.} Note that $\plim\, n^{-1}\SSR_j = \Var(r_j)$ (a constant), so $\SSR_j^2$ grows at rate $n^2$.

  \smallskip

  \textbf{Step 2.} For each $i$: $\hat{u}_i = u_i + (\beta - \hat{\beta})' x_i$. Under MLR.1--4 and consistency of $\hat{\beta}$:
  \[
    \hat{u}_i^2 \overset{p}{\approx} u_i^2 \quad \text{as } n \to \infty
  \]

  \smallskip

  \textbf{Step 3.} Therefore:
  \[
    \frac{1}{n}\sum_i \hat{r}_{ij}^2 \hat{u}_i^2
    \overset{p}{\to}
    \E[\hat{r}_{ij}^2 u_i^2] = \E[\hat{r}_{ij}^2 \sigma_i^2]
  \]

  \begin{keybox}
    The robust estimator substitutes $\hat{u}_i^2$ for $\sigma_i^2$ term-by-term. By the Law of Large Numbers, the average converges to the correct quantity. No assumption about the \textit{form} of $\sigma_i^2$ is required.
  \end{keybox}

\end{frame}

\begin{frame}{Robust $t$ and $F$ Statistics}

  Once we have robust SEs, all inference proceeds with the usual formulas.

  \textbf{Robust $t$-statistic:}
  \[
    t_{\text{rob}} = \frac{\hat{\beta}_j - \beta_j^0}{\widehat{\text{se}}_{\text{rob}}(\hat{\beta}_j)}
    \overset{d}{\to} \mathcal{N}(0,1) \quad \text{under } H_0, \; n \to \infty
  \]
  Same critical values as the standard $t$-test — only the SE changes.

  \smallskip

  \begin{highlightbox}
    \textbf{Robust $F$-statistic (Wald form)} — for $q$ restrictions $H_0: \mathbf{R}\boldsymbol{\beta} = \mathbf{r}$:
    \[
      F_{\text{rob}} = \frac{1}{q}(\mathbf{R}\hat{\boldsymbol{\beta}} - \mathbf{r})^{\prime}
      \bigl(\mathbf{R}\,\widehat{V}_{\text{rob}}\,\mathbf{R}^{\prime}\bigr)^{-1}
      (\mathbf{R}\hat{\boldsymbol{\beta}} - \mathbf{r})
      \overset{d}{\to} F_{q,\infty} \text{ under } H_0
    \]
    $\widehat{V}_{\text{rob}}$: robust variance-covariance matrix using White estimator block.
  \end{highlightbox}

\end{frame}

\begin{frame}{Example 8.1: OLS vs.\ Robust SEs — Log Wage Equation}

  \textbf{Dataset:} WAGE1 ($n=526$, $R^2=0.461$). Compare OLS SE (\ ) and robust SE [\ ].

  \begin{examplebox}{Log Wage Regression}
    \begin{center}
      \footnotesize
      \begin{tabular}{lcc}
        \toprule
        Variable & Coef. & OLS SE \quad [Robust SE] \\
        \midrule
        \textit{marrmale}  & $0.213$    & $(0.055)$ \quad $[0.057]$ \\
        \textit{marrfem}   & $-0.198$   & $(0.058)$ \quad $[0.058]$ \\
        \textit{singfem}   & $-0.110$   & $(0.056)$ \quad $[0.057]$ \\
        \textit{educ}      & $0.0789$   & $(0.0067)$ \quad $[0.0074]$ \\
        \textit{exper}     & $0.0268$   & $(0.0052)$ \quad $[0.0051]$ \\
        $\textit{exper}^2$ & $-0.00054$ & $(0.00011)$ \quad $[0.00011]$ \\
        \textit{tenure}    & $0.0291$   & $(0.0068)$ \quad $[0.0069]$ \\
        $\textit{tenure}^2$& $-0.00053$ & $(0.00023)$ \quad $[0.00024]$ \\
        intercept          & $0.321$    & $(0.100)$ \quad $[0.109]$ \\
        \bottomrule
      \end{tabular}
    \end{center}
  \end{examplebox}

\end{frame}

\begin{frame}{Example 8.1: Interpretation}

  \begin{keybox}
    \begin{itemize}
      \item Robust SE for \textit{educ}: $0.0067 \to 0.0074$ (slightly larger — mild heteroskedasticity present).
      \item Most SEs are nearly identical — significance conclusions unchanged.
      \item \textit{marrmale} premium: 21.3\% wage gap; \textit{marrfem} penalty: $-$19.8\%.
      \item Lesson: always report robust SEs and verify they do not alter inference.
    \end{itemize}
  \end{keybox}

  \smallskip

  \begin{highlightbox}
    When OLS and robust SEs are close, heteroskedasticity is \textbf{mild}. When they diverge substantially, inference from usual SEs is unreliable — robust SEs are preferred.
  \end{highlightbox}

\end{frame}

\begin{frame}[shrink=15]{Example 8.2: Robust $F$-Test for Race Dummies}

  \textbf{Dataset:} GPA3 ($n=366$). Dep.\ var.: \textit{cumgpa}.

  \begin{examplebox}{Cumulative GPA Regression}
    \begin{center}
      \footnotesize
      \begin{tabular}{lcc}
        \toprule
        Variable & Coef. & OLS SE \quad [Robust SE] \\
        \midrule
        \textit{sat}    & $0.00114$ & $(0.00018)$ \quad $[0.00019]$ \\
        \textit{hsperc} & $-0.00857$ & $(0.00124)$ \quad $[0.00140]$ \\
        \textit{tothrs} & $0.00250$ & $(0.00073)$ \quad $[0.00073]$ \\
        \textit{female} & $0.303$ & $(0.059)$ \quad $[0.059]$ \\
        \textit{black}  & $-0.128$ & $(0.147)$ \quad $[0.118]$ \\
        \textit{white}  & $-0.059$ & $(0.141)$ \quad $[0.110]$ \\
        intercept       & $1.47$ & $(0.23)$ \quad $[0.22]$ \\
        \midrule
        \multicolumn{3}{l}{$R^2 = 0.4006$} \\
        \bottomrule
      \end{tabular}
    \end{center}
  \end{examplebox}

  \begin{keybox}
    Testing $H_0: \beta_{\textit{black}} = \beta_{\textit{white}} = 0$: Usual $F = 0.40$ ($p=0.672$); Robust $F \approx 0.75$ — both fail to reject. Robust SEs for race dummies are \textit{smaller}.
  \end{keybox}

\end{frame}

\begin{frame}{Heteroskedasticity-Robust LM Test (\S8-2a)}

  \textbf{LM tests} can be made robust without matrix algebra.

  \begin{highlightbox}
    \textbf{Robust LM Test Procedure} (for testing $q$ exclusion restrictions):
    \begin{enumerate}
      \item Estimate the \textbf{restricted} model; save residuals $\tilde{u}_i$.
      \item Regress each \textbf{excluded} variable on all \textbf{included} variables; save residuals $\tilde{r}_{i\ell}$ for $\ell = 1,\ldots,q$.
      \item Regress \textbf{1} on $\tilde{r}_{i1}\tilde{u}_i, \ldots, \tilde{r}_{iq}\tilde{u}_i$ (no intercept); get $\SSR_0$.
      \item $\text{Robust LM} = n - \SSR_0 \overset{d}{\to} \chi^2_q$ under $H_0$.
    \end{enumerate}
  \end{highlightbox}

  \smallskip

  \begin{keybox}
    Key insight: the procedure uses products $\tilde{r}_{i\ell}\tilde{u}_i$ to ``orthogonalise'' the score. This makes the statistic valid under any form of heteroskedasticity — no variance model is assumed.
  \end{keybox}

\end{frame}

\begin{frame}{Example 8.3: Robust LM vs.\ Usual LM (CRIME1)}

  \textbf{Dataset:} CRIME1 ($n=2725$). Test $H_0: \beta_{\textit{avgsen}} = \beta_{\textit{avgsen}^2} = 0$.

  \begin{examplebox}{Arrest Rate Equation — LM Test Comparison}
    \begin{itemize}
      \item Usual $LM = 3.52$ ($p \approx 0.17$) — \textcolor{positive}{\textbf{fail to reject}} $H_0$
      \item Robust $LM = 5.10$ ($p \approx 0.078$) — \textcolor{negative}{\textbf{borderline reject}} at 10\%
    \end{itemize}
  \end{examplebox}

  \smallskip

  \begin{keybox}
    Heteroskedasticity \textit{distorts} the usual LM test. The robust version detects evidence against $H_0$ that the usual test misses. Using the wrong test can lead to incorrect conclusions about which variables belong in the model.
  \end{keybox}

\end{frame}

%% ============================================================
\section{Testing for Heteroskedasticity (\S8-3)}
%% ============================================================

\begin{frame}{Should We Always Use Robust SEs?}

  \begin{columns}[t]
    \column{0.49\textwidth}
    \textbf{Case for always-robust:}
    \begin{itemize}
      \item Valid under both homo- and heteroskedasticity
      \item No cost to unbiasedness
      \item Standard in modern applied work (cross-section data)
    \end{itemize}

    \column{0.49\textwidth}
    \textbf{Case for testing first:}
    \begin{itemize}
      \item Robust SEs can be imprecise in small samples
      \item Homoskedasticity $\Rightarrow$ OLS $t$-test is exact in finite samples (MLR.1--6)
      \item Knowing the variance form enables efficient WLS
    \end{itemize}
  \end{columns}

  \smallskip

  \begin{keybox}
    \textbf{Wooldridge's recommendation:}
    \begin{enumerate}
      \item \textbf{Always} report robust SEs in empirical work with cross-sectional data.
      \item Run a formal test if you want to understand the \textit{form} of heteroskedasticity — this guides whether to pursue WLS for efficiency gains.
    \end{enumerate}
  \end{keybox}

\end{frame}

\begin{frame}{The Breusch-Pagan Test: Setup}

  \textbf{Null hypothesis:} $H_0: \Var(u \mid x_1, \ldots, x_k) = \sigma^2$ (homoskedasticity)

  \smallskip

  \textbf{Idea:} Under $H_0$, $u_i^2$ should be unrelated to $x_1, \ldots, x_k$. Test this with a regression.

  \smallskip

  \begin{definitionbox}{BP Test Auxiliary Regression (eq.\ 8.12)}
    \[
      \hat{u}_i^2 = \delta_0 + \delta_1 x_{i1} + \delta_2 x_{i2} + \cdots + \delta_k x_{ik} + v_i
    \]
    Joint null: $H_0: \delta_1 = \delta_2 = \cdots = \delta_k = 0$ (eq.\ 8.13)
  \end{definitionbox}

  \smallskip

  \textbf{Why squared residuals?} $\hat{u}_i^2$ is the sample proxy for $\sigma_i^2 = \Var(u_i \mid x_i)$. If the latter depends on $x$, so should $\hat{u}_i^2$.

\end{frame}

\begin{frame}{The Breusch-Pagan Test: Test Statistics}

  \begin{highlightbox}
    \textbf{Two equivalent test statistics:}

    \smallskip

    \textbf{F-statistic} from the auxiliary regression [8.12]:
    \[
      F = \frac{R_{\hat{u}^2}^2 / k}{(1 - R_{\hat{u}^2}^2)/(n-k-1)} \sim F_{k,\,n-k-1} \quad \text{under } H_0
    \]

    \smallskip

    \textbf{LM statistic} (Lagrange Multiplier):
    \[
      LM = n \cdot R_{\hat{u}^2}^2 \overset{d}{\to} \chi^2_k \quad \text{under } H_0 \quad \text{(eq.\ 8.16)}
    \]
    where $R_{\hat{u}^2}^2$ is the $R^2$ from regressing $\hat{u}_i^2$ on $x_1, \ldots, x_k$.
  \end{highlightbox}

  \smallskip

  \begin{keybox}
    \textbf{Limitation:} BP detects only \textit{linear} dependence of $\sigma_i^2$ on $x_j$.

    If variance depends nonlinearly on $x_j$ (e.g., $\sigma_i^2 = \exp(\delta_0 + \delta_1 x_{i1}^2)$), BP may fail to reject even when heteroskedasticity is present. $\Rightarrow$ Use the White test.
  \end{keybox}

\end{frame}

\begin{frame}{Example 8.4: BP Test — Levels vs.\ Logs}

  \begin{columns}[t]
    \column{0.48\textwidth}
    \begin{examplebox}{Level Model (HPRICE1)}
      \begin{align*}
        \widehat{\textit{price}} &= -21.77 + 0.00207\,\textit{lotsize} \\
                                 &\quad + 0.123\,\textit{sqrft} + 13.85\,\textit{bdrms}
      \end{align*}
      $n=88$, $R^2=0.672$

      \smallskip
      BP $F = 5.34$, $p = 0.002$\\
      BP $LM = 14.09$, $p = 0.003$\\
      \smallskip
      \textcolor{negative}{\textbf{Reject} $H_0$} — heteroskedastic.
    \end{examplebox}

    \column{0.48\textwidth}
    \begin{examplebox}{Log-Log Model (HPRICE1)}
      \begin{align*}
        \widehat{\log(\textit{price})} &= -1.30 + 0.168\log(\textit{lotsize}) \\
                                       &\quad + 0.700\log(\textit{sqrft}) + 0.037\,\textit{bdrms}
      \end{align*}
      $n=88$, $R^2=0.643$

      \smallskip
      BP $F = 1.41$, $p = 0.245$\\
      BP $LM = 4.22$, $p = 0.239$\\
      \smallskip
      \textcolor{positive}{\textbf{Fail to reject}} $H_0$.
    \end{examplebox}
  \end{columns}

  \smallskip

  \begin{keybox}
    Log-transforming the dependent variable \textit{stabilises the variance}. Level models for prices, wages, and income frequently exhibit heteroskedasticity that disappears in log form.
  \end{keybox}

\end{frame}

\begin{frame}{The White Test: Motivation}

  \textbf{Problem with BP:} Misses heteroskedasticity that depends on \textit{squares and cross-products} of regressors.

  \smallskip

  \textbf{White (1980) proposed} a more general auxiliary regression:

  \begin{definitionbox}{White Test — Full Form (eq.\ 8.19)}
    With regressors $x_1, x_2, \ldots, x_k$, regress $\hat{u}_i^2$ on all unique terms:
    \[
      x_1, x_2, \ldots, x_k, \quad x_1^2, x_2^2, \ldots, x_k^2, \quad x_1 x_2, x_1 x_3, \ldots, x_{k-1}x_k
    \]
    Joint $F$-test or $LM = n R^2_{\hat{u}^2} \overset{d}{\to} \chi^2_m$ where $m$ = number of slope terms.
  \end{definitionbox}

  \smallskip

  \begin{keybox}
    \textbf{Cost of the full White test:} With $k$ regressors, there are $k + k + \binom{k}{2} = k(k+3)/2$ terms. For $k=5$, that is 20 regressors in the auxiliary regression — many degrees of freedom lost.
  \end{keybox}

\end{frame}

\begin{frame}{The White Test: Parsimonious Special Case}

  \textbf{Insight:} The fitted value $\hat{y}_i = \hat{\beta}_0 + \hat{\beta}_1 x_{i1} + \cdots + \hat{\beta}_k x_{ik}$ is a linear combination of all $x_j$, so $\hat{y}_i^2$ is a linear combination of all cross-products.

  \smallskip

  \begin{definitionbox}{White Test — Special Case (eq.\ 8.20)}
    Regress $\hat{u}_i^2$ on the OLS fitted values and their squares:
    \[
      \hat{u}_i^2 = \delta_0 + \delta_1 \hat{y}_i + \delta_2 \hat{y}_i^2 + \text{error}
    \]
    $H_0: \delta_1 = \delta_2 = 0$. \quad
    $LM = n R^2_{\hat{u}^2} \overset{d}{\to} \chi^2_2$ — only \textbf{2 degrees of freedom} regardless of $k$.
  \end{definitionbox}

  \smallskip

  \begin{keybox}
    \textbf{Advantage:} Degrees of freedom do not grow with $k$. Powerful for detecting common heteroskedasticity patterns (variance increasing with $\hat{y}$).

    \smallskip

    \textbf{Disadvantage:} May miss heteroskedasticity that is unrelated to the direction of $\hat{y}$.
  \end{keybox}

\end{frame}

\begin{frame}{Example 8.5: White Test on the Log Housing Price Equation}

  \textbf{Recall:} BP test on the log-log model failed to reject homoskedasticity ($p = 0.239$).

  \smallskip

  \begin{examplebox}{Special White Test}
    Regress $\hat{u}_i^2$ on $\widehat{\log(\textit{price})}_i$ and $\widehat{\log(\textit{price})}_i^2$:
    \[
      R_{\hat{u}^2}^2 = 0.0392
    \]
    \[
      LM = 88 \times 0.0392 = 3.45, \quad p\text{-value} = 0.178 \quad (2 \text{ d.f.})
    \]
  \end{examplebox}

  \smallskip

  \begin{itemize}
    \item Fails to reject at 10\% level (slightly more evidence than BP, but inconclusive).
    \item Both tests agree: the log-log model is approximately homoskedastic.
  \end{itemize}

  \smallskip

  \begin{keybox}
    \textbf{Practical conclusion:} Use the log-log specification for housing prices. Robust SEs are still advisable as a precaution, but WLS is likely unnecessary here.
  \end{keybox}

\end{frame}

%% ============================================================
\section{Weighted Least Squares (\S8-4)}
%% ============================================================

\begin{frame}{The Logic of WLS: Re-weighting Observations}

  \textbf{Motivation:} If we know how variance varies, we can do better than OLS.

  \smallskip

  \begin{definitionbox}{WLS Setup (eq.\ 8.21)}
    Suppose the heteroskedasticity has the form:
    \[
      \Var(u_i \mid x_i) = \sigma^2 h_i, \quad h_i = h(x_{i1}, \ldots, x_{ik}) > 0
    \]
    where $h_i$ is known (or estimable).
  \end{definitionbox}

  \smallskip

  \begin{highlightbox}
    \textbf{WLS minimises the weighted sum of squared residuals (eq.\ 8.27):}
    \[
      \sum_{i=1}^n \frac{(y_i - \beta_0 - \beta_1 x_{i1} - \cdots - \beta_k x_{ik})^2}{h_i}
    \]
    This \textbf{down-weights} noisy observations ($h_i$ large) and \textbf{up-weights} precise ones ($h_i$ small).
  \end{highlightbox}

\end{frame}

\begin{frame}{WLS as OLS on a Transformed Model}

  \textbf{Divide the regression equation through by $\sqrt{h_i}$:}

  Starting from (eq.\ 8.24): $y_i = \beta_0 + \beta_1 x_{i1} + \cdots + \beta_k x_{ik} + u_i$

  \smallskip

  Dividing through gives the \textbf{transformed model} (eq.\ 8.25--8.26):
  \[
    y_i^* = \beta_0 x_{i0}^* + \beta_1 x_{i1}^* + \cdots + \beta_k x_{ik}^* + u_i^*
  \]
  where $y_i^* = y_i/\sqrt{h_i}$, \; $x_{ij}^* = x_{ij}/\sqrt{h_i}$, \; $x_{i0}^* = 1/\sqrt{h_i}$, \; $u_i^* = u_i/\sqrt{h_i}$.

  \smallskip

  \begin{highlightbox}
    \textbf{Proof that transformed model is homoskedastic:}
    \[
      \Var(u_i^* \mid x_i) = \frac{\Var(u_i \mid x_i)}{h_i} = \frac{\sigma^2 h_i}{h_i} = \sigma^2
      \quad \checkmark
    \]
    Since MLR.5 now holds for the transformed model, \textbf{OLS on the starred variables is BLUE} (Gauss-Markov).
  \end{highlightbox}

\end{frame}

\begin{frame}{WLS: Key Properties and Warning}

  \begin{highlightbox}
    \textbf{If $h_i$ is correctly specified:}
    \begin{itemize}
      \item WLS is \textbf{unbiased} (linear transformation of the data)
      \item WLS is \textbf{BLUE} — lower variance than OLS
      \item WLS standard errors (from the transformed regression) are valid
    \end{itemize}
  \end{highlightbox}

  \smallskip

  \begin{keybox}
    \textbf{Warning — misspecification risk:}

    If $h_i$ is \textit{wrong}, WLS is still consistent but:
    \begin{itemize}
      \item WLS standard errors may be \textbf{invalid}
      \item WLS may even be \textbf{less efficient} than OLS
    \end{itemize}
    \textbf{Safe practice:} After WLS, always compute \textit{heteroskedasticity-robust SEs} for the WLS estimates. These are valid even if $h_i$ is misspecified.
  \end{keybox}

\end{frame}

\begin{frame}{Example 8.6: Financial Wealth with WLS}

  \textbf{Setup:} Assume $\Var(u_i \mid \textit{inc}_i) = \sigma^2 \cdot \textit{inc}_i$, so $h_i = \textit{inc}_i$. Dataset: 401KSUBS.

  \begin{examplebox}{Financial Wealth Equations (SEs in parentheses)}
    \begin{center}
      \footnotesize
      \begin{tabular}{lcc}
        \toprule
        Variable & OLS & WLS ($w_i = 1/\textit{inc}_i$) \\
        \midrule
        \textit{inc} (simple) & $0.821\;(0.104)$ & $0.787\;(0.063)$ \\
        intercept             & $-10.57\;(2.53)$ & $-9.58\;(1.65)$ \\
        \midrule
        \textit{inc} (expanded) & $0.771\;(0.100)$ & $0.740\;(0.064)$ \\
        $(\textit{age}-25)^2$   & $0.0251\;(0.0043)$ & $0.0175\;(0.0019)$ \\
        \textit{male}           & $2.48\;(2.06)$ & $1.84\;(1.56)$ \\
        \textit{e401k}          & $6.89\;(2.29)$ & $5.19\;(1.70)$ \\
        \bottomrule
      \end{tabular}
    \end{center}
  \end{examplebox}

\end{frame}

\begin{frame}{Example 8.6: Interpretation — WLS vs.\ OLS}

  \begin{keybox}
    \begin{itemize}
      \item WLS SEs are \textbf{uniformly smaller} — efficiency gain from correct weighting.
      \item 401(k) eligibility (\textit{e401k}): coefficient drops from 6.89 (OLS) to 5.19 (WLS) — \textit{economically meaningful} (roughly \$1,700 wealth difference).
      \item WLS relies on the assumption $\Var(u \mid \textit{inc}) \propto \textit{inc}$. If wrong, efficiency advantage disappears.
    \end{itemize}
  \end{keybox}

  \smallskip

  \begin{highlightbox}
    \textbf{Caution:} If the variance function $h_i$ is misspecified, always report \textbf{robust WLS SEs} rather than conventional WLS SEs to protect against invalid inference.
  \end{highlightbox}

\end{frame}

\begin{frame}{Feasible GLS: When $h_i$ Must Be Estimated (\S8-4b)}

  In practice, the variance function $h(x)$ is almost never known exactly.

  \smallskip

  \begin{definitionbox}{FGLS: Exponential Variance Function (eq.\ 8.30)}
    Assume:
    \[
      \Var(u_i \mid x_i) = \sigma^2 \exp(\delta_0 + \delta_1 x_{i1} + \cdots + \delta_k x_{ik})
    \]
    Taking logs of $u_i^2 \approx \sigma^2 \exp(\delta_0 + \delta_1 x_{i1} + \cdots + \delta_k x_{ik}) \cdot \epsilon_i$:
    \[
      \log(\hat{u}_i^2) = \alpha_0 + \delta_1 x_{i1} + \cdots + \delta_k x_{ik} + \text{error}
      \quad \text{(eq.\ 8.32)}
    \]
    Estimated weights: $\hat{h}_i = \exp(\hat{g}_i)$ where $\hat{g}_i$ are fitted values from [8.32].
  \end{definitionbox}

\end{frame}

\begin{frame}{Feasible GLS: Step-by-Step Algorithm}

  \begin{highlightbox}
    \textbf{FGLS Procedure:}
    \begin{enumerate}
      \item Estimate main equation by \textbf{OLS}; obtain residuals $\hat{u}_i$.
      \item Compute $\log(\hat{u}_i^2)$.
      \item Regress $\log(\hat{u}_i^2)$ on $x_{i1}, \ldots, x_{ik}$; save fitted values $\hat{g}_i$.
      \item Form estimated weights: $\hat{h}_i = \exp(\hat{g}_i)$.
      \item Apply \textbf{WLS} with weights $w_i = 1/\hat{h}_i$.
      \item Report \textbf{robust SEs} for the WLS step (in case $h_i$ is misspecified).
    \end{enumerate}
  \end{highlightbox}

  \smallskip

  \begin{keybox}
    \textbf{Properties of FGLS:}
    \begin{itemize}
      \item \textbf{Consistent and asymptotically more efficient} than OLS (if form is correct)
      \item \textbf{Not unbiased} — $\hat{h}_i$ introduces additional estimation uncertainty
      \item With misspecified $h_i$: still consistent, but efficiency not guaranteed
    \end{itemize}
  \end{keybox}

\end{frame}

\begin{frame}{Example 8.7: Cigarette Demand — OLS vs.\ FGLS}

  \textbf{Dataset:} SMOKE ($n=807$). BP test: $LM=32.3$ ($p\approx 0$) — strong heteroskedasticity confirmed.

  \begin{examplebox}{Cigarette Demand: OLS vs. FGLS (SEs in parentheses)}
    \begin{center}
      \footnotesize
      \begin{tabular}{lcc}
        \toprule
        Variable & OLS & FGLS \\
        \midrule
        $\log(\textit{income})$   & $0.880\;(0.728)$  & $1.30\;(0.44)$ \\
        $\log(\textit{cigprice})$ & $-0.751\;(5.773)$ & $-2.94\;(4.46)$ \\
        \textit{educ}             & $-0.501\;(0.167)$ & $-0.463\;(0.120)$ \\
        \textit{age}              & $0.771\;(0.160)$  & $0.482\;(0.097)$ \\
        $\textit{age}^2$          & $-0.0090\;(0.0017)$ & $-0.0056\;(0.0009)$ \\
        \textit{restaurn}         & $-2.83\;(1.11)$   & $-3.46\;(0.80)$ \\
        \midrule
        $R^2$ & 0.053 & 0.113 \\
        \bottomrule
      \end{tabular}
    \end{center}
  \end{examplebox}

\end{frame}

\begin{frame}{Example 8.7: Interpretation — FGLS Efficiency Gains}

  \begin{keybox}
    \begin{itemize}
      \item FGLS SEs are \textbf{uniformly smaller} — substantial efficiency gain from re-weighting.
      \item Income elasticity: $0.880 \to 1.30$ (FGLS more precise and estimate shifts — OLS was noisy).
      \item Restaurant smoking restrictions: $-3.46$ cigs/day under FGLS vs.\ $-2.83$ OLS — larger and more precise.
      \item $R^2$ nearly doubles: $0.053 \to 0.113$ — FGLS fits the weighted data much better.
    \end{itemize}
  \end{keybox}

  \smallskip

  \begin{highlightbox}
    When heteroskedasticity is \textbf{strong} (BP $p \approx 0$), FGLS can dramatically improve efficiency and even shift point estimates. But remember: FGLS requires correct specification of $h(x)$.
  \end{highlightbox}

\end{frame}

\begin{frame}{Estimation Strategy Summary}

  \begin{highlightbox}
    \textbf{Deciding which estimator and standard errors to use:}
    \begin{center}
      \footnotesize
      \begin{tabular}{lll}
        \toprule
        \textbf{Situation} & \textbf{Estimator} & \textbf{Standard Errors} \\
        \midrule
        Homoskedasticity confirmed & OLS & Usual OLS SEs \\
        Heterosk., form unknown & OLS & White robust SEs \\
        Heterosk., form known & WLS & Usual WLS SEs \\
        Heterosk., form estimated & FGLS & Robust WLS SEs \\
        \bottomrule
      \end{tabular}
    \end{center}
  \end{highlightbox}

  \smallskip

  \begin{keybox}
    When in doubt: \textbf{OLS + robust SEs} is the defensible baseline. Pursue WLS/FGLS only when you have a strong theoretical or empirical reason to believe a specific variance function.
  \end{keybox}

\end{frame}

%% ============================================================
\section{The Linear Probability Model Revisited (\S8-5)}
%% ============================================================

\begin{frame}{LPM Heteroskedasticity: The Theoretical Result}

  Recall from Lecture 6: the LPM has \textit{built-in} heteroskedasticity.

  \begin{highlightbox}
    \textbf{Proof:} For the LPM with $y_i \in \{0,1\}$, given $x_i$, $y_i$ is Bernoulli with mean $p_i = \E[y_i \mid x_i]$.

    \smallskip
    Therefore: $\Var(y_i \mid x_i) = p_i(1-p_i) = p(x_i)[1-p(x_i)]$ \quad (eq.\ 8.45)

    \smallskip
    Since $p(x_i) = \beta_0 + \beta_1 x_{i1} + \cdots + \beta_k x_{ik}$ depends on $x_i$, so does the variance.

    MLR.5 is \textbf{always violated} in the LPM.
  \end{highlightbox}

  \smallskip

  \begin{keybox}
    \textbf{Natural WLS weights for LPM} (eq.\ 8.47):
    \[
      \hat{h}_i = \hat{y}_i(1 - \hat{y}_i)
    \]
    \textbf{Constraint:} WLS is only feasible if all $\hat{y}_i \in (0,1)$. If any $\hat{y}_i \leq 0$ or $\hat{y}_i \geq 1$, weights are non-positive and WLS fails.
  \end{keybox}

\end{frame}

\begin{frame}{Example 8.8: Married Women's Labor Force Participation}

  \textbf{Dataset:} MROZ ($n=753$). $y = \textit{inlf}$ (1 = in labor force).

  \begin{examplebox}{LPM for Labor Force Participation}
    \begin{center}
      \footnotesize
      \begin{tabular}{lcc}
        \toprule
        Variable & Coef. & OLS SE \quad [Robust SE] \\
        \midrule
        \textit{nwifeinc} & $-0.0034$ & $(0.0014)$ \quad $[0.0015]$ \\
        \textit{educ}     & $0.038$   & $(0.007)$ \quad $[0.007]$ \\
        \textit{exper}    & $0.039$   & $(0.006)$ \quad $[0.006]$ \\
        $\textit{exper}^2$& $-0.00060$ & $(0.00012)$ \quad $[0.00012]$ \\
        \textit{age}      & $-0.016$  & $(0.002)$ \quad $[0.002]$ \\
        \textit{kidslt6}  & $-0.262$  & $(0.034)$ \quad $[0.034]$ \\
        \textit{kidsge6}  & $0.013$   & $(0.013)$ \quad $[0.013]$ \\
        intercept         & $0.586$   & $(0.154)$ \quad $[0.151]$ \\
        \midrule
        \multicolumn{3}{l}{$n=753$, $R^2=0.264$} \\
        \bottomrule
      \end{tabular}
    \end{center}
  \end{examplebox}

  \begin{keybox}
    \begin{itemize}
      \item OLS and robust SEs are very similar — heteroskedasticity is mild here.
      \item Each young child (\textit{kidslt6}) reduces LFP probability by \textbf{26.2 pp} (large!).
      \item Each year of education raises LFP by \textbf{3.8 pp}.
      \item Conclusion: OLS + robust SEs is sufficient; WLS offers little gain.
    \end{itemize}
  \end{keybox}

\end{frame}

\begin{frame}{Example 8.9: PC Ownership — OLS vs.\ WLS for LPM}

  \textbf{Dataset:} GPA1 ($n=141$). $y = \textit{PC}$ (1 = owns computer).

  \begin{examplebox}{PC Ownership: OLS vs. WLS}
    \begin{center}
      \footnotesize
      \begin{tabular}{lcc}
        \toprule
        Variable & OLS & WLS ($w_i = 1/\hat{h}_i$) \\
        \midrule
        \textit{hsGPA}  & $0.065\;(0.137)$   & $0.033\;(0.130)$ \\
        \textit{ACT}    & $0.00065\;(0.0155)$ & $0.0043\;(0.0155)$ \\
        \textit{parcoll}& $0.221\;(0.093)$   & $0.215\;(0.086)$ \\
        intercept       & $-0.0004\;(0.490)$ & $0.026\;(0.477)$ \\
        \midrule
        $R^2$ & 0.0415 & 0.0464 \\
        \bottomrule
      \end{tabular}
    \end{center}
  \end{examplebox}

  \begin{keybox}
    \begin{itemize}
      \item OLS and WLS give very similar results — no efficiency concern here.
      \item \textit{parcoll} is the only significant predictor: parental college degree raises PC ownership by $\approx 22$ pp.
      \item With mild heteroskedasticity and small $n$, OLS + robust SEs is preferred.
    \end{itemize}
  \end{keybox}

\end{frame}

%% ============================================================
\section{Summary and Decision Framework}
%% ============================================================

\begin{frame}[shrink=8]{Three-Step Decision Framework}

  \begin{highlightbox}
    \textbf{Step 1 — Estimate by OLS; report with robust SEs.}

    OLS is unbiased and consistent under MLR.1--4. Robust SEs ensure valid inference regardless of $\sigma_i^2$.

    \smallskip

    \textbf{Step 2 — Test for heteroskedasticity (optional but informative).}

    Run BP test ($\chi^2_k$) and/or White special case ($\chi^2_2$). If rejected, proceed to Step 3.

    \smallskip

    \textbf{Step 3 — Pursue efficiency if heteroskedasticity confirmed.}
    \begin{itemize}
      \item \textit{Known} $h_i$: use WLS with conventional WLS SEs.
      \item \textit{Estimated} $h_i$: use FGLS with \textbf{robust} WLS SEs.
      \item \textit{Unknown form}: stay with OLS + robust SEs.
    \end{itemize}
  \end{highlightbox}

\end{frame}

\begin{frame}[shrink=8]{Key Concepts Reviewed}

  \begin{keybox}
    \textbf{The Problem:} Heteroskedasticity ($\sigma_i^2 \neq \sigma^2$) leaves $\hat{\beta}_j$ unbiased but invalidates SEs and tests (MLR.5 violated).

    \smallskip
    \textbf{Fix for Inference:} White robust SE — $\widehat{\Var}(\hat{\beta}_j) = \sum_i \hat{r}_{ij}^2 \hat{u}_i^2 / \SSR_j^2$ — consistent under any $\sigma_i^2$.

    \smallskip
    \textbf{Testing:} BP test ($\hat{u}^2$ on $x_j$, $\chi^2_k$) or White special case ($\hat{u}^2$ on $\hat{y}, \hat{y}^2$, $\chi^2_2$).

    \smallskip
    \textbf{Fix for Efficiency:} WLS/FGLS achieves BLUE when $h_i = \Var(u_i \mid x_i)/\sigma^2$ is correctly specified. Use robust WLS SEs when $h_i$ is estimated.
  \end{keybox}

\end{frame}

\begin{frame}{Bridge to Lecture 8: Specification and Data Issues}

  \begin{columns}[t]
    \column{0.48\textwidth}
    \begin{highlightbox}
      \textbf{Lecture 7 (today)} \\
      MLR.5 violated only:
      \begin{itemize}
        \item OLS is still unbiased and consistent
        \item Only SEs and tests are affected
        \item Fix: robust SEs or WLS
      \end{itemize}
    \end{highlightbox}

    \column{0.48\textwidth}
    \begin{highlightbox}
      \textbf{Lecture 8 (Chapter 9)} \\
      MLR.4 violated ($\E[u \mid x] \neq 0$):
      \begin{itemize}
        \item Omitted variable bias
        \item Functional form misspecification
        \item Measurement error (attenuation)
        \item Missing data / selection bias
      \end{itemize}
    \end{highlightbox}
  \end{columns}

  \smallskip

  \begin{keybox}
    \small
    \begin{tabular}{lll}
      \toprule
      Lecture & Assumption violated & Effect on OLS \\
      \midrule
      7 (today) & MLR.5 only & Biased SEs; $\hat{\beta}$ consistent \\
      8 (next)  & MLR.4      & $\hat{\beta}$ biased and inconsistent \\
      \bottomrule
    \end{tabular}
  \end{keybox}

\end{frame}

%% ============================================================
\end{document}
